% Part: normal-modal-logic
% Chapter: tableaux
% Section: completeness

\documentclass[../../../include/open-logic-section]{subfiles}

\begin{document}

\olfileid[fr]{nml}{tab}{cpl}

\olsection{Complétude pour \Log{K}}

\begin{explain}
  Pour montrer que la méthode des !!{tableau}s est complète, nous
  devons établir que, s'il n'existe aucun !!{tableau} fermé démontrant
  $\Gamma \Proves !A$, alors $\Gamma \Entails/ !A$ : il existe un
  contre-modèle. L'absence de !!{tableau} fermé signifie que toute
  tentative de construction laisse au moins une branche ouverte.
  L'idée est d'adopter une procédure \emph{systématique et exhaustive}
  d'application des règles. Si un !!{tableau} fermé existe, cette
  procédure finit aussi par fermer toutes ses branches. Si elle ne
  ferme pas toutes les branches,
  une branche ouverte complète fournit les données nécessaires à un
  contre-modèle. Les mondes de ce modèle seront tous les préfixes de
  cette branche, et !!a{propositional variable}~$p$ sera vraie au monde
  $\sigma$ si et seulement si $\sFmla{\True}{p}[\sigma]$ y figure.
\end{explain}

\begin{defn}
  Une branche d'!!a{tableau} est dite complète si elle est fermée ou
  si, chaque fois qu'elle contient une !!{formula} préfixée
  $\sFmla{S}{!A}[\sigma]$ à laquelle une règle peut s'appliquer, elle
  contient aussi :
  \begin{enumerate}
    \item les !!{formula}s préfixées qui sont les conclusions de la
      règle, pour les règles propositionnelles sans bifurcation ;
    \item toutes les !!{formula}s conclusives d’au moins l’une des branches produites, pour
      les règles propositionnelles avec bifurcation ;
    \item au moins une conclusion possible, pour les règles modales
      qui exigent un nouveau préfixe ;
    \item la conclusion correspondante pour chaque préfixe successeur
      de la forme $\sigma.n$ déjà utilisé sur la branche, pour les
      règles modales qui exigent un tel préfixe.
  \end{enumerate}
\end{defn}

\begin{explain}
Par exemple, une branche complète contient $\sFmla{\True}{!B}[\sigma]$
et $\sFmla{\True}{!C}[\sigma]$ chaque fois qu'elle contient
$\sFmla{\True}{!B \land !C}[\sigma]$. Si elle contient
$\sFmla{\True}{!B \lor !C}[\sigma]$, elle contient au moins l'une
des formules $\sFmla{\True}{!B}[\sigma]$ et
$\sFmla{\True}{!C}[\sigma]$. Si elle contient \iftag{prvBox}
{$\sFmla{\False}{\Box !B}[\sigma]$} {$\sFmla{\True}{\Diamond !B}[\sigma]$},
elle contient également \iftag{prvBox}
{$\sFmla{\False}{!B}[\sigma.n]$}
{$\sFmla{\True}{!B}[\sigma.n]$} pour au moins un~$n$. Enfin,
si elle contient \iftag{prvBox} {$\sFmla{\True}{\Box !B}[\sigma]$}
{$\sFmla{\False}{\Diamond !B}[\sigma]$}, elle contient également \iftag{prvBox}
{$\sFmla{\True}{!B}[\sigma.n]$}
{$\sFmla{\False}{!B}[\sigma.n]$} pour chaque~$n$ tel que
$\sigma.n$ est déjà utilisé sur la branche.
\end{explain}

\begin{prop}\ollabel{prop:complete-tableau}
  Tout ensemble fini $\Gamma$ admet !!a{tableau} dont chaque branche
  est complète.
\end{prop}

\begin{proof}
  Considérons une branche ouverte d'!!a{tableau} pour~$\Gamma$. À
  chaque étape, elle contient un nombre fini de !!{formula}s préfixées
  auxquelles une règle peut s'appliquer. Dans un ordre fixé, par
  exemple de haut en bas, examinons chacune de ces !!{formula}s
  préfixées. Pour chaque !!{formula} pour laquelle les conditions
  (1)--(4) ne sont pas déjà remplies, appliquons les règles
  pertinentes. Lorsqu'une règle scinde la branche, nous
  poursuivons sur chacune des deux branches résultantes. Nous
  répétons l'opération sur toute branche encore ouverte et
  incomplète.

  Si $\Gamma$ est vide, aucune règle ne s'applique. Sinon,
  soit $d(!A)$ la profondeur modale de~$!A$, et soit $D$
  le maximum de $|\sigma|+d(!A)$ pour les formules signées
  $\sFmla{S}{!A}[\sigma]$ de~$\Gamma$, où $|\sigma|$ désigne
  la longueur du préfixe. Les règles propositionnelles conservent
  le préfixe et ne font pas croître la profondeur modale. Les
  règles modales passent de $\sigma$ à $\sigma.n$ en supprimant
  une modalité extérieure. La somme reste donc au plus égale
  à~$D$.

  Sur une branche, les formules à chaque préfixe proviennent de
  l'ensemble fini des sous-formules signées de~$\Gamma$. Chacune
  des formules qui exigent un témoin ne peut créer qu'un
  successeur à ce préfixe, car le premier témoin satisfait
  déjà la condition (3). Aux successeurs ainsi créés s'ajoutent
  seulement ceux qui sont présents dès le départ. Le nombre de
  successeurs immédiats est donc uniformément borné, et la
  longueur des préfixes est bornée par~$D$. Il ne peut y avoir
  qu'un nombre fini de préfixes sur une branche. Chaque
  prolongement ajoute au moins une formule signée à chacun de
  ses résultats, et il n'y a qu'un nombre fini de ces formules
  à chaque préfixe. La construction se termine. À la fin, toute branche est
  fermée ou satisfait les conditions (1)--(4) : elle est donc
  complète.
\end{proof}

\begin{thm}[Complétude]
  \ollabel{thm:tableau-completeness}
  Si $\Gamma$ n'admet aucun !!{tableau} fermé, $\Gamma$ est
  satisfaisable.
\end{thm}

\begin{proof}
Si $\Gamma$ est vide, un modèle à un monde et l’interprétation vide le satisfont. Supposons donc $\Gamma$ fini et non vide. D'après la proposition, il admet
!!a{tableau} dont chaque branche est complète. Puisqu'aucun
!!{tableau} pour~$\Gamma$ n'est fermé, il existe !!a{tableau} avec
au moins une branche ouverte et complète. Soit $\Delta$ l'ensemble
des !!{formula}s préfixées de cette branche, et $P(\Delta)$ l'ensemble
des préfixes qui y figurent.

Définissons le modèle~$\mModel{M(\Delta)} = \tuple{P(\Delta), R, V}$.
Ses mondes sont les préfixes de~$\Delta$, sa relation d'accessibilité
est définie par
\[
R\sigma\sigma' \quad \text{si et seulement si} \quad
\sigma'=\sigma.n \quad \text{pour un certain } n
\]
et
\[
V(p) = \Setabs{\sigma}{\sFmla{\True}{p}[\sigma] \in \Delta}.
\]
Montrons par récurrence sur~$!A$ que, si $\sFmla{\True}{!A}[\sigma] \in
\Delta$, alors $\mSat{M(\Delta)}{!A}[\sigma]$, et que, si
$\sFmla{\False}{!A}[\sigma] \in \Delta$, alors
$\mSat/{M(\Delta)}{!A}[\sigma]$.

\begin{enumerate} \iftag{prvFalse}{\item \indcase{!A}{\lfalse}{Le signe vrai fermerait la branche ; le signe faux est satisfait dans tout monde.}}{} \iftag{prvTrue}{\item \indcase{!A}{\ltrue}{Le signe faux fermerait la branche ; le signe vrai est satisfait dans tout monde.}}{} \iftag{prvIff}{\item \indcase{!A}{!B\liff!C}{Par complétude de la branche, le signe vrai donne soit les deux opérandes vrais, soit les deux opérandes faux ; le signe faux donne des signes opposés. L’hypothèse de récurrence sur les deux opérandes établit exactement la valeur du biconditionnel.}}{}
  \item \indcase{!A}{p}{Si $\sFmla{\True}{\indfrm}[\sigma] \in \Delta$,
    alors $\sigma \in V(p)$ par définition de~$V$, d'où
    $\mSat{M(\Delta)}{\indfrm}[\sigma]$.

    Si $\sFmla{\False}{\indfrm}[\sigma] \in \Delta$, alors
    $\sFmla{\True}{\indfrm}[\sigma] \notin \Delta$, car la branche
    serait autrement fermée. Ainsi, $\sigma \notin V(p)$, et donc
    $\mSat/{M(\Delta)}{\indfrm}[\sigma]$.}
  \item \indcase{!A}{\lnot !B}{\iftag{probNot}{Exercice.}{Si
      $\sFmla{\True}{\indfrm}[\sigma] \in \Delta$, alors
      $\sFmla{\False}{!B}[\sigma] \in \Delta$, puisque la branche est
      complète. Par hypothèse de récurrence,
      $\mSat/{M(\Delta)}{!B}[\sigma]$, et donc
      $\mSat{M(\Delta)}{\indfrm}[\sigma]$.

      Si $\sFmla{\False}{\indfrm}[\sigma] \in \Delta$, alors
      $\sFmla{\True}{!B}[\sigma] \in \Delta$ puisque la branche est
      complète. Par hypothèse de récurrence,
      $\mSat{M(\Delta)}{!B}[\sigma]$, et donc
      $\mSat/{M(\Delta)}{\indfrm}[\sigma]$.}}
  \item \indcase{!A}{!B \land !C}{\iftag{probAnd}{Exercice.}{Si
      $\sFmla{\True}{\indfrm}[\sigma] \in \Delta$, alors
      $\sFmla{\True}{!B}[\sigma] \in \Delta$ et
      $\sFmla{\True}{!C}[\sigma] \in \Delta$, puisque la branche est
      complète. Par hypothèse de récurrence,
      $\mSat{M(\Delta)}{!B}[\sigma]$ et
      $\mSat{M(\Delta)}{!C}[\sigma]$. Donc
      $\mSat{M(\Delta)}{\indfrm}[\sigma]$.

      Si $\sFmla{\False}{\indfrm}[\sigma] \in \Delta$, alors
      $\sFmla{\False}{!B}[\sigma] \in \Delta$ ou
      $\sFmla{\False}{!C}[\sigma] \in \Delta$, puisque la branche est
      complète. Par hypothèse de récurrence,
      $\mSat/{M(\Delta)}{!B}[\sigma]$ ou
      $\mSat/{M(\Delta)}{!C}[\sigma]$. Donc
      $\mSat/{M(\Delta)}{\indfrm}[\sigma]$.}}
  \item \indcase{!A}{!B \lor !C}{\iftag{probOr}{Exercice.}{Si
      $\sFmla{\True}{\indfrm}[\sigma] \in \Delta$, alors
      $\sFmla{\True}{!B}[\sigma] \in \Delta$ ou
      $\sFmla{\True}{!C}[\sigma] \in \Delta$, puisque la branche est
      complète. Par hypothèse de récurrence,
      $\mSat{M(\Delta)}{!B}[\sigma]$ ou
      $\mSat{M(\Delta)}{!C}[\sigma]$. Donc
      $\mSat{M(\Delta)}{\indfrm}[\sigma]$.

      Si $\sFmla{\False}{\indfrm}[\sigma] \in \Delta$, alors
      $\sFmla{\False}{!B}[\sigma] \in \Delta$ et
      $\sFmla{\False}{!C}[\sigma] \in \Delta$, puisque la branche est
      complète. Par hypothèse de récurrence,
      $\mSat/{M(\Delta)}{!B}[\sigma]$ et
      $\mSat/{M(\Delta)}{!C}[\sigma]$. Donc
      $\mSat/{M(\Delta)}{\indfrm}[\sigma]$.}}
  \item \indcase{!A}{!B \lif !C}{\iftag{probIf}{Exercice.}{Si
      $\sFmla{\True}{\indfrm}[\sigma] \in \Delta$, alors
      $\sFmla{\False}{!B}[\sigma] \in \Delta$ ou
      $\sFmla{\True}{!C}[\sigma] \in \Delta$, puisque la branche est
      complète. Par hypothèse de récurrence,
      $\mSat/{M(\Delta)}{!B}[\sigma]$ ou
      $\mSat{M(\Delta)}{!C}[\sigma]$. Donc
      $\mSat{M(\Delta)}{\indfrm}[\sigma]$.

      Si $\sFmla{\False}{\indfrm}[\sigma] \in \Delta$, alors
      $\sFmla{\True}{!B}[\sigma] \in \Delta$ et
      $\sFmla{\False}{!C}[\sigma] \in \Delta$, puisque la branche est
      complète. Par hypothèse de récurrence,
      $\mSat{M(\Delta)}{!B}[\sigma]$ et
      $\mSat/{M(\Delta)}{!C}[\sigma]$. Donc
      $\mSat/{M(\Delta)}{\indfrm}[\sigma]$.}}
  \item \indcase{!A}{\Box !B}{\iftag{probBox}{Exercice.}{Si
      $\sFmla{\True}{\indfrm}[\sigma] \in \Delta$, alors, comme la
      branche est complète, $\sFmla{\True}{!B}[\sigma.n] \in \Delta$
      pour chaque $\sigma.n$ déjà utilisé sur la branche, c'est-à-dire
      pour chaque $\sigma' \in P(\Delta)$ tel que $R\sigma\sigma'$.
      Par hypothèse de récurrence, $\mSat{M(\Delta)}{!B}[\sigma']$
      pour chaque $\sigma'$ tel que $R\sigma\sigma'$. Par conséquent,
      $\mSat{M(\Delta)}{\indfrm}[\sigma]$.

      Si $\sFmla{\False}{\indfrm}[\sigma] \in \Delta$, alors il existe
      un $\sigma.n$ tel que $\sFmla{\False}{!B}[\sigma.n] \in \Delta$,
      puisque la branche est complète. Par hypothèse de récurrence,
      $\mSat/{M(\Delta)}{!B}[\sigma.n]$. Comme $R\sigma(\sigma.n)$,
      il existe un~$\sigma'$ tel que
      $\mSat/{M(\Delta)}{!B}[\sigma']$. Donc
      $\mSat/{M(\Delta)}{\indfrm}[\sigma]$.}}
  \item \indcase{!A}{\Diamond !B}{\iftag{probDiamond}{Exercice.}{Si
      $\sFmla{\True}{\indfrm}[\sigma] \in \Delta$, alors il existe
      un $\sigma.n$ tel que $\sFmla{\True}{!B}[\sigma.n] \in \Delta$,
      puisque la branche est complète. Par hypothèse de récurrence,
      $\mSat{M(\Delta)}{!B}[\sigma.n]$. Comme $R\sigma(\sigma.n)$,
      il existe un~$\sigma'$ tel que
      $\mSat{M(\Delta)}{!B}[\sigma']$. Donc
      $\mSat{M(\Delta)}{\indfrm}[\sigma]$.

      Si $\sFmla{\False}{\indfrm}[\sigma] \in \Delta$, alors, comme la
      branche est complète, $\sFmla{\False}{!B}[\sigma.n] \in \Delta$
      pour chaque $\sigma.n$ déjà utilisé sur la branche, c'est-à-dire
      pour chaque $\sigma' \in P(\Delta)$ tel que $R\sigma\sigma'$.
      Par hypothèse de récurrence,
      $\mSat/{M(\Delta)}{!B}[\sigma']$ pour chaque
      $\sigma'$ tel que $R\sigma\sigma'$. Par conséquent,
      $\mSat/{M(\Delta)}{\indfrm}[\sigma]$.}}
\end{enumerate}
L'identité sur~$P(\Delta)$ est une interprétation des préfixes dans
$\mModel{M(\Delta)}$. Comme $\Gamma \subseteq \Delta$,
$\mSat{M(\Delta)}{\Gamma}[\mathrm{id}]$.

Si $\Gamma$ est infini et n'admet aucun tableau fermé, toute
partie finie de~$\Gamma$ est satisfaisable d'après le cas fini.
Dans le langage de la traduction standard, introduisons une
constante $c_\sigma$ pour chaque $\sigma\in P(\Gamma)$. À
chaque formule signée au préfixe $\sigma$, associons la
traduction de la formule au monde $c_\sigma$, précédée d'une
négation si le signe est $\False$. Ajoutons
$\Atom{Q}{c_\sigma,c_{\sigma.n}}$ chaque fois que les deux
préfixes appartiennent à~$P(\Gamma)$, où $Q$ désigne la relation
d'accessibilité.

Considérons une partie finie de cet ensemble d'énoncés. Dans
$\Gamma$, prenons les formules signées dont elle contient la
traduction et, pour chaque préfixe nommé dans une contrainte
d'accessibilité de cette partie, une formule de~$\Gamma$ qui
porte ce préfixe. Cet ensemble fini est satisfaisable.
L'interprétation de ses préfixes fournit la valeur des
constantes considérées et vérifie les contraintes
d'accessibilité entre elles. On obtient ainsi un modèle de la
partie finie. Par compacité du premier ordre, l'ensemble des
énoncés admet un modèle. Les interprétations des constantes
$c_\sigma$ y donnent une interprétation des préfixes qui
satisfait~$\Gamma$.
\end{proof}

\begin{prob}
Complétez la démonstration du \olref[nml][tab][cpl]{thm:tableau-completeness}.
\end{prob}

\begin{cor}
\ollabel{cor:entailment-completeness}
Si $\Gamma \Entails !A$, alors $\Gamma \Proves !A$.
\end{cor}

\begin{cor}
\ollabel{cor:weak-completeness}
Si $!A$ est vraie dans tous les modèles, alors $\Proves !A$.
\end{cor}
\begin{explain}
La saturation des branches fournit une procédure finie lorsque
les hypothèses sont en nombre fini. Pour un ensemble infini,
la compacité donne l'existence d'un modèle en réunissant dans
une même interprétation les contraintes exprimées par les
formules signées et par les relations entre préfixes.
\end{explain}
\end{document}
